\documentclass[11pt,a4paper]{article}

\usepackage[margin=2.5cm]{geometry}
\usepackage[T1]{fontenc}
\usepackage{lmodern}
\usepackage{microtype}
\usepackage{amsmath,amssymb}
\usepackage{booktabs,tabularx,array}
\usepackage{enumitem}
\usepackage{needspace}
\usepackage[numbers,sort&compress]{natbib}
\usepackage{xcolor}
\usepackage[colorlinks=true,allcolors=blue!60!black]{hyperref}

\title{Perishable Inventory-Routing under Synthetic Demand Uncertainty: An Operational Comparison of Forecast-Based and Decision-Focused Policies}
\author{}
\date{}

\begin{document}
\maketitle

\begin{abstract}
Perishable inventory-routing couples replenishment, age-specific carryover, expiry, customer service, and shared vehicle capacity. We study these trade-offs in a five-period, ten-store synthetic benchmark and compare point-forecast and quantile predict-then-optimize (PTO), a disclosed adaptation of a published deliver-up-to policy, and an approximate two-phase decision-focused learning (DFL) hybrid. Demand is hand-generated; no M5 observations or operational data are used. Across 60 independent held-out episodes, DFL had mean cost 2795.0 model units versus 2738.1 for point PTO; the paired difference was 56.9 (95\% episode-bootstrap interval $[-14.2,129.1]$, Holm-adjusted $p=0.137$). DFL was lower-cost than the open-loop UL$_\lambda$-inspired heuristic adaptation, but that is not a faithful reproduction of the published policy. A post-hoc finite-scenario robust adaptation had mean cost 2712.3, but all 300 of its candidate-plan MILP calls used the construction fallback. All 300 primary evaluation plans from the forecast-based/DFL methods also used fallback, and no evaluation MILP was optimal. Results are limited to this synthetic setting and do not establish field performance or state-of-the-art superiority; no source-faithful recent robust or scenario-based operational comparator is evaluated.
\end{abstract}

\noindent\textbf{Keywords:} perishable inventory; inventory-routing; age-specific inventory; demand uncertainty; fleet capacity; predict-then-optimize; decision-focused learning

\section{Introduction}

A distributor serving retail stores must decide when to replenish, how much product to deliver, and how to route a limited vehicle fleet. With perishable goods, these choices also determine the age profile of store inventory, the amount that expires, and the stock available to serve later demand. Inventory carried today can reduce a future trip or shortage, but its value depends on remaining shelf life and storage capacity. A method that forecasts demand accurately may still produce poor operating decisions if forecast errors interact with route-capacity thresholds, expiry, or asymmetric shortage costs.

This study takes those inventory and routing decisions as its subject. We ask: \emph{under a controlled synthetic demand process, how do forecast-based and decision-focused methods compare in operating cost, service, expiry, carryover, and routing when they share the same physical inventory evaluator?} Decision-focused learning (DFL) is treated as one candidate approach, alongside conventional point and quantile PTO and an age-aware replenishment-policy adaptation. We do not assume DFL is superior.

Perishable inventory theory has long studied age and outdating \citep{nahmias1982perishable,haijema2014new}. Deterministic perishable inventory-routing research develops exact formulations and branch-and-cut methods \citep{coelho2014optimal,alvarez2020formulations}; stochastic perishable inventory-routing work studies demand uncertainty and compares ordering and decomposition methods \citep{soysal2015modeling,crama2018stochastic}. More recent work studies robust perishable PIRP \citep{he2024robust}. In adjacent online inventory routing, Jia et al. \citep{jia2025scenario} propose quantile-based scenario prediction followed by scenario optimization. These studies establish that uncertainty-aware perishable inventory-routing and online scenario methods already exist. Our computational comparisons are narrower and do not reproduce those methods.

The contribution of this paper is therefore limited to a transparent, reproducible synthetic comparison with explicit physical outcomes and information timing. The benchmark separates expiry spoilage from capacity-driven disposal, measures unmet demand and fill rate, tracks age-specific on-hand inventory, and records route and solver diagnostics. The methods and evaluator are described so that the limits of each comparison are visible. The present experiment is a single base configuration; it does not estimate general shelf-life, cost-ratio, or network-topology effects.

\section{Related literature and comparator scope}

\subsection{Perishable inventory and inventory-routing}

Perishable inventory models represent finite product life, age-dependent stock, issuing rules, and outdating \citep{nahmias1982perishable,haijema2014new}. In a deterministic PIRP, replenishment and vehicle routes are jointly selected over time; exact formulations and hybrid heuristics have been developed for fixed-life products \citep{coelho2014optimal,alvarez2020formulations}. Under stochastic demand, Crama et al. \citep{crama2018stochastic} compare an expected-value policy, the UL$_\lambda$ deliver-up-to policy, and stochastic decomposition and decomposition-integration methods. Their UL$_\lambda$ policy is dynamic: it updates the store state after observing each period's demand. He et al. \citep{he2024robust} study robust perishable PIRP and propose robust-counterpart reformulation, Benders decomposition, and an ILS-BD heuristic.

These are direct precedents for the physical problem and uncertainty, and they define a higher comparison bar than forecast-only ablations. Our UL$_\lambda$-inspired code does not reproduce the source policy: it fixes an open-loop five-period plan, uses validation residuals and predicted state propagation, and constructs routes with a nearest-neighbor/2-opt heuristic. Its results are attributed only to this adaptation. The study also includes a post-hoc finite-scenario robust plan-selection adaptation motivated by the direct robust perishable-PIRP problem class. It does not implement He et al.'s uncertainty-set robust counterpart, Benders decomposition, or ILS-BD, and it is reported separately as exploratory evidence.

Jia et al. \citep{jia2025scenario} provide a recent data-driven online inventory-routing comparator. Their full paper describes multi-horizon quantile-RNN forecasts, sampled demand scenarios, and scenario-based stochastic optimization. The formulation reviewed has no finite-life age-cohort state; adding FIFO age tracking and expiry would be an adaptation. Its online decision timing also differs from the episode-start fixed-plan protocol used here, so a fair comparison requires a common information schedule. No such implementation is included in the present run. Deterministic, non-perishable multivehicle IRP matheuristics are also relevant algorithmically but would require adaptation before serving as perishable stochastic comparators \citep{lagana2024ils}.

\subsection{Predict-then-optimize and decision-focused learning}

PTO estimates demand and then supplies forecasts to an operational optimizer. We compare a point predictor and a quantile predictor to distinguish forecast loss choices while holding the downstream planning model fixed. DFL trains a predictive model against a surrogate of downstream decision quality; SPO+ is a standard surrogate framework for linear optimization with uncertain objective coefficients \citep{elmachtoub2022smart}. Surveys and benchmarks cover DFL methods across a range of combinatorial problems \citep{mandi2024decision}. However, an approximate oracle or a custom surrogate does not inherit the guarantees of exact SPO+ with an exact optimization oracle. We accordingly describe our implemented learner as an approximate DFL hybrid and state its oracle use and fallbacks.

\section{System, demand process, and information timing}

\subsection{Perishable inventory-routing setting}

The benchmark contains one depot, ten stores, two homogeneous vehicles, and a five-period planning horizon. Store coordinates are selected from a clustered Solomon C1 layout; Euclidean distances are used, and Solomon time windows are not modeled \citep{solomon1987algorithms}. Each vehicle carries at most 100 units per period; each store holds at most 70 units. Product shelf life is three periods. Starting inventory is zero at every store. Within the horizon, inventory ages by one period, demand consumes oldest stock first, unconsumed stock at expiry is discarded, and excess surviving inventory above store capacity is removed oldest-first.

Let $I_{ita}$ denote store $i$'s end-of-period inventory of age $a$, $q_{it}$ the delivered quantity, and $d_{it}$ demand in period $t$. The physical evaluator applies FIFO to available opening inventory and the fresh delivery, records unmet demand $u_{it}$, records expiry $e_{it}$ separately from capacity disposal $c_{it}$, and carries forward only surviving age cohorts. The planning model jointly represents age-structured inventory, delivery quantities, and vehicle routes. Forecast-based methods pass a five-period forecast into that model. All planned schedules are then evaluated against the same realized demand using the FIFO simulator.

Per-episode operating cost is
\begin{equation}
 C = C^{\mathrm{route}} + C^{\mathrm{hold}} + p_s\sum_{i,t}(e_{it}+c_{it}) + p_u\sum_{i,t}u_{it},
 \label{eq:cost}
\end{equation}
where route cost is total distance, holding cost is age-weighted on-hand inventory, and $p_s=5$, $p_u=20$, and base holding cost $h=1$. We report expiry and capacity disposal separately even though both receive penalty $p_s$. Fill rate is $1-\sum_{i,t}u_{it}/\sum_{i,t}d_{it}$.

\subsection{Synthetic demand and sample design}

Demand is generated by \texttt{data/generate\_dataset.py} using illustrative store-specific means, day-of-week seasonality, annual trend, announced promotions and discounts, a price response, hidden synthetic weather effects, and autoregressive noise. This is a hand-specified simulation; parameters are not estimated from M5 or operational records. Solomon coordinates describe benchmark geometry only and do not constitute an observed distribution network.

The final dataset contains 400 independently seeded episodes, each with a private 28-day history warmup and a five-day decision horizon. Episodes are split into 280 training, 60 validation, and 60 test cases. Predictors receive store identity, calendar variables, pre-episode demand summaries, and announced promotion/price inputs. Realized within-horizon demand and realized future weather are excluded. All methods receive the same episode-start feature set, and test outcomes are not used for training or tuning. The independent unit for inference is the held-out episode; the three DFL training seeds are averaged within each episode.

The 60-episode design was selected under a normal-approximation target of 80\% power for a paired standardized effect of $d=0.45$ at two-sided $\alpha=0.025$ for two primary comparisons. This is a design assumption, not a guarantee of achieved power or a substitute for reporting uncertainty.

\section{Methods}

\subsection{Forecast-based policies}

\textbf{LightGBM point PTO} fits a 100-tree LightGBM regressor (maximum depth 5; learning rate 0.05) to training episodes. Nonnegative five-period predictions are sent to the common age-structured planning model.

\textbf{LightGBM quantile PTO} uses the same learner settings and fits the inventory-cost critical quantile $\alpha=p_u/(p_u+h+p_s/\tau)$. Its predictions are passed to the same planning model. It is a cost-sensitive forecasting baseline, not a published PIRP method.

\subsection{UL$_\lambda$-inspired adaptation}

We form store-period residuals from the validation episodes and tune $\lambda\in\{1,2,3\}$ and target service level in $\{0.80,0.90,0.95,0.98\}$ by validation operating cost. For each planned period, the policy uses residual samples to estimate whether current predicted inventory covers one-period demand; if not, it targets a validation quantile of cumulative predicted demand over up to $\lambda$ periods, subject to store capacity. Target deliveries are packed into vehicle-capacity-feasible routes by nearest neighbor and 2-opt. Inventory for later periods is propagated under predicted demand, not observed test demand. Thus it is a fixed-horizon adaptation inspired by Crama et al.'s published dynamic UL$_\lambda$ rule, not a reproduction or a source-faithful rolling-horizon baseline.

\subsection{Exploratory finite-scenario robust plan selection}

After the revision2 pilot outcomes had been inspected, we added an exploratory robust-policy adaptation. Validation-only residual paths form an empirical uncertainty set that preserves observed cross-store and cross-period residual dependence. For each test episode, the procedure solves five deterministic age-structured PIRPs using demand paths at prespecified aggregate-residual quantiles $0,0.25,0.50,0.75,1.00$. It evaluates each candidate's single fixed delivery-and-route schedule against every path in the validation residual set with the common FIFO cost, then selects the candidate with the smallest worst-case cost. Test demand is not used to generate scenarios or select a plan. This is a restricted finite-scenario candidate-selection heuristic motivated by robust PIRP work \citep{he2024robust}; it is not a reproduction of He et al.'s robust-counterpart reformulation, Benders decomposition, or ILS-BD. Because it was added after the pilot outcomes were inspected, its results are exploratory. All five candidate MILP requests per episode returned the feasibility-checked construction fallback in the evaluated run.

\subsection{Approximate two-phase DFL hybrid}

The learner is a neural demand model trained with a two-phase schedule and an asymmetric pinball component. It uses a continuous LP warmup, followed by an integer-oracle phase. For the 50-epoch run, the implementation activates a fast path: cached LP decisions and a vectorized capacity-projected polyhedral loss are used on most minibatches; the configured HiGHS integer oracle is requested on selected minibatches in the final phase. The model's SPO+-inspired loss also relies on solver decisions that can be heuristic fallbacks. Therefore the method is not a full-MILP-per-minibatch SPO+ implementation, and no exact-oracle SPO+ guarantee is claimed. Checkpoints are selected by realized validation operating cost, not by regret against a heuristic reference.

\subsection{Optimization status and fallback}

Forecast-based plans and selected DFL oracle calls request SciPy/HiGHS MILPs with relative MIP gap 0.01 and a 0.5-second time limit. The solver output is checked against the model constraints. If the solver returns no acceptable incumbent, or if a feasible construction has lower cost than a non-optimal incumbent, the feasibility-checked LP-guided construction is used. Each forecast plan is checked against the FIFO simulator under its own planning demand; a maximum absolute state discrepancy above $10^{-4}$ stops the run. We report requested settings, actual solver status/gap, incumbent/fallback use, and parity. “MILP requested” is not reported as “MILP optimal.”

\section{Evaluation and statistical analysis}

The primary outcome is total operating cost per five-period test episode. Physical secondary outcomes are routing cost/distance, holding cost, expiry cost and units, capacity-disposal cost and units, stockout cost and unmet units, demand-weighted fill rate, delivered units, average and ending on-hand inventory, age-specific carryover, number of routes, and fleet-capacity utilization. Forecast MAE/RMSE are diagnostic measures only.

The prespecified primary pairwise comparisons are DFL minus LightGBM point PTO and DFL minus the UL$_\lambda$-inspired adaptation. Differences are computed on common test episodes. We report the paired mean difference, paired standard deviation, 95\% episode-bootstrap interval (10,000 replicates), raw paired $t$-test $p$-value, and Holm-adjusted $p$-value over the two primary comparisons. Quantile PTO and disaggregated physical measures are secondary/descriptive. Training seeds are averaged within episode, not treated as extra independent cases. No regret against a heuristic full-information plan is used as the primary metric.

\section{Results}

The run contains 60 independent test episodes for each method; DFL results are averaged across three training seeds within each episode. Table~\ref{tab:primary-results} reports episode-level means and standard deviations for cost, with physical quantities averaged over episodes. Cost values are model units defined by the synthetic routing and penalty parameters.

\begin{table}[htbp]
\centering
\small
\caption{Primary held-out outcomes (60 independent episodes). Total cost is mean (SD); other columns are episode means.}
\label{tab:primary-results}
\begin{tabular}{lrrrrr}
\toprule
Method & Total cost (SD) & Fill rate & Unmet units & Expired units & Capacity discards \\
\midrule
UL$_\lambda$ adap. & 3058.3 (1282.5) & 0.886 & 121.6 & 0.00 & 0.00 \\
Point PTO & 2738.1 (1176.6) & 0.901 & 106.2 & 0.73 & 0.00 \\
Quantile PTO & 3238.9 (1393.9) & 0.885 & 123.9 & 1.05 & 0.00 \\
DFL hybrid & 2795.0 (1212.1) & 0.903 & 104.1 & 0.83 & 0.00 \\
\bottomrule
\end{tabular}
\end{table}

\begin{table}[htbp]
\centering
\small
\caption{Mean cost components by method (model units per episode). Capacity disposal was zero in all methods.}
\label{tab:cost-components}
\begin{tabular}{lrrrrr}
\toprule
Method & Routing & Holding & Expiry & Stockout & Total \\
\midrule
UL$_\lambda$ adap. & 483.4 & 143.4 & 0.0 & 2431.4 & 3058.3 \\
Point PTO & 465.5 & 144.0 & 3.7 & 2124.9 & 2738.1 \\
Quantile PTO & 464.3 & 290.8 & 5.3 & 2478.5 & 3238.9 \\
DFL hybrid & 467.3 & 241.5 & 4.2 & 2082.0 & 2795.0 \\
\bottomrule
\end{tabular}
\end{table}

The prespecified paired cost comparison found no statistically detectable cost reduction for DFL relative to point PTO: DFL minus point PTO was 56.9 model units (95\% paired-bootstrap interval $[-14.2,129.1]$, paired $t$-test $p=0.137$, Holm-adjusted $p=0.137$). DFL was lower-cost than the UL$_\lambda$ adaptation by 263.3 units (95\% interval $[-340.0,-183.4]$, Holm-adjusted $p=2.75\times10^{-8}$). That comparison is against this study's open-loop nearest-neighbor/2-opt adaptation, not Crama et al.'s dynamic UL$_\lambda$ method or their decomposition methods. Quantile PTO and the age-specific physical measures are descriptive; no additional significance claims are made.

We separately evaluated the post-hoc finite-scenario robust adaptation on the same 60 test episodes. It had mean cost 2712.3 (SD 1148.9), mean fill rate 0.904, mean expired quantity 0.64 units, and zero capacity discards.

\Needspace{5\baselineskip}
Its exploratory paired cost difference was $-25.8$ versus point PTO (95\% episode-bootstrap interval $[-85.2,32.7]$) and $-82.7$ versus DFL (interval $[-150.0,-14.9]$). These contrasts were not prespecified and do not establish the performance of He et al.'s source method. All 300 candidate-plan MILP requests selected the construction fallback; the comparison therefore evaluates this finite-scenario heuristic with fallback plans, not optimal robust solutions.

Mean delivered quantities were 942.3, 947.3, 991.4, and 991.6 units for the UL$_\lambda$ adaptation, point PTO, quantile PTO, and DFL, respectively. Mean on-hand inventory per store-period was 2.87, 2.88, 5.81, and 4.83 units, while mean ending inventory was 40.1, 29.1, 90.5, and 71.1 units. Mean age-1/age-2 on-hand stock per store-period was 2.86/0.01, 2.69/0.18, 5.64/0.18, and 4.67/0.15 units, respectively. Each method used 10 routes per five-period episode on average; mean fleet-capacity utilization ranged from 0.942 to 0.992. These physical summaries describe this generated demand and the specified cost parameters only.

Cost accounting reconciled exactly, the data hashes matched the run manifest, all 60 held-out IDs were covered by the four primary methods, and maximum plan/FIFO state discrepancy was zero. The separate robust adaptation also covered all 60 test episodes and had zero maximum candidate-plan/FIFO state discrepancy. The solver audit is a material limitation: all 300 primary evaluation plan calls and all 300 robust candidate-plan requests selected the feasibility-checked construction fallback; no evaluation MILP was proven optimal. DFL training made 1,534 integer-oracle requests, all of which fell back; the other 936 training solver calls were LP calls. These results compare implemented policies under a shared physical evaluator, not exact optimization methods.

\Needspace{9\baselineskip}
\section{Discussion and limitations}

Interpretation will be limited to the stated ten-store C1 geometry, five-period horizon, zero initial inventory, one shelf life, one cost setting, hand-specified synthetic demand process, and solver/time-budget configuration. The test episodes provide independent simulation replications, not a sample of real retail networks. The simulation does not validate observed expiry, customer service, or route execution. Its demand features and cost parameters are illustrative rather than empirically calibrated.

The operational comparator set remains narrow. The UL$_\lambda$ adaptation is open-loop and uses a route-packing heuristic. The finite-scenario robust adaptation is a custom, post-hoc method with a restricted candidate set; it is not He et al.'s source algorithm. DFL did not show lower cost than point PTO in the primary paired comparison, and all evaluation MILPs fell back to feasible constructions. Jia et al.'s online scenario PTO and He et al.'s robust perishable PIRP are real and relevant methods but neither source method has been reproduced. The study therefore cannot establish state-of-the-art performance or resolve the editor's comparator concern. A subsequent study should reproduce a contemporary robust or scenario operational method with an adequately functioning solver and a matched information protocol. Broader shelf-life/capacity testing and operational validation would further strengthen the paper.

\section{Conclusion}

This paper reframes the problem around the physical perishable inventory-routing decision and reports a controlled synthetic comparison with common FIFO accounting. Its claims are restricted to the implemented policies, tested configuration, and solver outcomes. The current evidence can characterize operating trade-offs in that benchmark; it cannot support claims of empirical field performance or state-of-the-art superiority. Any stronger conclusion depends on closing the operational-comparator and validation gaps above.

\bibliographystyle{plainnat}
\bibliography{paper/references}
\end{document}
